\subsection{初始拓扑}

命题 4. 设 X 是一个集合，设 \((\mathbf{Y}_i)_{i \in I}\) 是拓扑空间的一个族，并且对每个 \(i \in I\)，设 \(f_i\) 是由 X 到 \(\mathbf{Y}_i\) 中的一个映射。设 \(\mathfrak{S}\) 是 X 的形如 \(\overline{f}_i^1 (\mathbf{U}_i)\)（\(i \in I\)，\(\mathbf{U}_i\) 在 \(\mathbf{Y}_i\) 中开）的子集所成之集，并设 \(\mathfrak{B}\) 是 \(\mathfrak{S}\) 的集合的有限交所成之集。则 \(\mathfrak{B}\) 是 X 上某个拓扑 \(\mathfrak{C}\) 的拓扑基，该拓扑是 X 关于族 \((f_i)\) 的初始拓扑结构（《集合论》，第四章，§ 2，第 3 目），特别地，它是使诸映射 \(f_i\) 连续的 X 上最粗的拓扑。更精确地说，若 \(g\) 是由拓扑空间 Z 到 X 中的映射，则 \(g\) 在某点 \(z \in Z\) 处连续（X 赋以拓扑 \(\mathfrak{C}\)）当且仅当每个函数 \(f_i \circ g\) 都在 \(z\) 处连续。

设 \(\mathfrak{D}\) 是属于 \(\mathfrak{B}\) 的集合的所有并所成之集；显然 \(\mathfrak{D}\) 满足公理 \((\mathrm{O}_{\mathrm{I}})\)，因为取并的运算是可结合的；而 \(\mathfrak{D}\) 满足公理 \((\mathrm{O}_{\mathrm{II}})\)，这是由于 \(\mathfrak{B}\) 的定义以及有限交对任意并的分配律 {[}Set Theory, R § 4, formula (37){]}。于是集合 \(\mathfrak{D}\) 是 X 关于某个拓扑 \(\mathfrak{G}\) 的开子集所成之集，\(\mathfrak{B}\) 是该拓扑的一个拓扑基。我们将证明命题的最后一个断言，由初始结构的一般性质（《集合论》，第四章，§ 2，第 3 目，判别准则 CST 9），它蕴含其余断言。首先，\(\mathfrak{S}\) 的定义表明诸 \(f_{i}\) 在 X 上连续（第 1 目，定理 1）；因此，若 \(g\) 在 \(z\) 处连续，则诸映射 \(f_{i} \circ g\) 亦然（第 1 目，命题 2）。反之，设所有映射 \(f_{i} \circ g\) 都在 \(z\) 处连续，并设 V 是 \(g(z)\) 在 X 中的一个邻域；按定义，存在 I 的一个有限子集 J，并且对每个 \(i \in J\)，存在开子集 \(\mathrm{U}_{i}\)（\(\mathrm{Y}_{i}\) 的开子集），使得 V 包含集合 \(\bigcap_{i \in J} f_{i}^{1}(\mathrm{U}_{i})\) 且 \(g(z)\) 属于这个集合。由此推出

\[
\overline {{g}} ^ {1} (\mathrm{V}) \supset \bigcap_ {i \in J} \overline {{g}} ^ {1} (\overline {{f}} _ {i} ^ {1} (\mathrm{U} _ {i})),
\]

而假设蕴含诸集合 \(\overline{g}^{1}(\overline{f}_{i}^{1}(\mathbf{U}_{i}))\) 中的每一个都是 \(z\) 在 \(\mathbf{Z}\) 中的邻域；于是 \(\overline{g}^{1}(\mathbf{V})\) 也是 \(z\) 在 \(\mathbf{Z}\) 中的邻域。证毕。

设 \(\mathfrak{B}_i\) 是 \(\mathbf{Y}_i(t\in I)\) 的拓扑的一个拓扑基；用 \(\mathfrak{S}'\) 表示 \(X\) 的形如 \(\overline{f}_i^1 (U_i)\) 的子集所成之集（对 \(t\in I\) 且 \(U_{t}\in \mathfrak{B}_{t}\)（对每个 \(t\in I\)））；若 \(\mathfrak{B}'\) 是 \(\mathfrak{S}'\) 的集合的有限交所成之集，则显然 \(\mathfrak{B}'\) 是拓扑 \(\mathfrak{C}\) 的一个拓扑基。

初始结构的一般性质（《集合论》，第四章，§ 2，第 3 目，判别准则 CST 10）特别地蕴含下列传递性质（其直接证明是十分容易的）：

命题 5. 设 X 是一个集合，\((Z_{t})_{t\in I}\) 是拓扑空间的一个族，\((J_{\lambda})_{\lambda\in L}\) 是 I 的一个分拆，而 \((Y_{\lambda})_{\lambda\in L}\) 是以 L 为指标的一族集合。又对每个 \(\lambda\in L\)，设 \(h_{\lambda}\) 是由 X 到 \(Y_{\lambda}\) 中的一个映射；对每个 \(\lambda\in L\) 与每个 \(t\in J_{\lambda}\)，设 \(g_{t\lambda}\) 是由 \(Y_{\lambda}\) 到 \(Z_{t}\) 中的一个映射，并令 \(f_{t}=g_{t\lambda}\circ h_{\lambda}\)。如果每个 \(Y_{\lambda}\) 都赋以使诸映射 \(g_{t\lambda}(t\in J_{\lambda})\) 连续的最粗拓扑，那么使诸 \(f_{t}\) 连续的 X 上最粗的拓扑与使诸 \(h_{\lambda}\) 连续的最粗拓扑相同。

例。1) 拓扑的逆像。设 X 是一个集合，Y 是拓扑空间，f 是由 X 到 Y 中的映射；使 f 连续的 X 上最粗的拓扑 C 称为 Y 的拓扑在 f 下的逆像。由命题 4 以及并与交的逆像公式 {[}Set Theory, R, § 4, formulae (34) and (46){]} 可知，拓扑 C 中的开（相应地，闭）集是 Y 的开（相应地，闭）集在 f 下的逆像；从而，对每个 \(x \in X\)，诸集合 \(\overline{f}^{1}(W)\)（其中 W 取遍 \(f(x)\) 在 Y 中的某个基本邻域系）构成 x 在拓扑 C 中的基本邻域系。在 § 3 中我们将以诱导拓扑之名研究 X 是 Y 的子集且 f 是典范单射 \(X \to Y\) 这一特殊情形；这时赋以诱导拓扑的 X 称为 Y 的子空间。

为使映射 \(f\)（由拓扑空间 \(\mathbf{X}\) 到拓扑空间 \(\mathbf{X}'\) 中的映射）连续，必要且充分的是 \(\mathbf{X}\) 的拓扑比 \(f\) 之下 \(\mathbf{X}'\) 的拓扑的逆像细。

\begin{enumerate}
\def\labelenumi{\arabic{enumi})}
\setcounter{enumi}{1}
\tightlist
\item
  拓扑集的上确界。拓扑的每个族 \((\mathcal{C}_t)_{i\in I}\)（集合 \(X\) 上的拓扑的族）都有一个上确界 \(\mathcal{C}\)（在 \(X\) 上所有拓扑的有序集中），即存在 \(X\) 上的一个拓扑，它在 \(X\) 上比每个 \(\mathcal{C}_t\) 都细的所有拓扑中是最粗的。为看出这一点，可以应用命题 4，取 \(Y_t\) 为集合 \(X\)（赋以拓扑 \(\mathcal{C}_t\)），取 \(f_t\) 为恒等映射 \(X\to Y_t\)；\(\mathcal{C}\) 是使所有映射 \(f_t\) 连续的最粗拓扑。
\end{enumerate}

设 \(\mathfrak{S}\) 是集合 X 的子集的任意集合；在 X 上的诸拓扑 \(\mathfrak{C}\)（使 \(\mathfrak{S}\) 的集合皆为开集者）中，有一个拓扑 \(\mathfrak{C}_0\) 比其他所有拓扑都粗，称为由 \(\mathfrak{S}\) 生成的拓扑。对每个集合 \(U \in \mathfrak{S}\)，设 \(\mathfrak{C}_U\) 为以 \(\varnothing\)、U 与 X 为开集的拓扑{[}显然 X 的这个子集集合满足 \((O_I)\) 与 \((O_{II})]\)；则 \(\mathfrak{C}_0\) 恰是诸拓扑 \(\mathfrak{C}_U\) 的上确界。由命题 4，若 \(\mathfrak{B}\) 是属于 \(\mathfrak{S}\) 的集合的有限交所成之集，则 \(\mathfrak{B}\) 是拓扑 \(\mathfrak{C}_0\) 的一个拓扑基。我们说 \(\mathfrak{S}\) 是 \(\mathfrak{C}_0\) 的一个子基。

\begin{enumerate}
\def\labelenumi{\arabic{enumi})}
\setcounter{enumi}{2}
\tightlist
\item
  积拓扑。设 \((\mathbf{X}_t)_{i \in I}\) 是拓扑空间的一个族。积集合 \(\mathbf{X} = \prod_{i \in I} \mathbf{X}_i\) 上使诸
\end{enumerate}

投影 \(pr_{i}: X \rightarrow X_{i}\) 为连续映射的最粗拓扑称为诸 \(X_{i}\) 的拓扑之积；我们将在 § 4 中更详细地研究它。

